\begin{table}[htbp]
\centering
\caption{Configurações utilizadas nos experimentos e respectivos tempos médios de inferência.}
\label{tab:configuracoes-experimentos-wasm}

\small
\begin{tabularx}{\textwidth}{
>{\raggedright\arraybackslash}p{2.3cm}
>{\raggedright\arraybackslash}X
>{\centering\arraybackslash}p{1.4cm}
>{\centering\arraybackslash}p{1.6cm}
>{\centering\arraybackslash}p{1.5cm}
>{\centering\arraybackslash}p{1.8cm}
}
\toprule
\textbf{Configuração} &
\textbf{Plataforma / ambiente} &
\textbf{CPU} &
\textbf{Memória} &
\textbf{Imagens} &
\textbf{Tempo médio} \\
\midrule

WebAssembly interpretado &
ESP32-CAM / WAMR &
160 MHz &
4 MB PSRAM &
6 &
144,019 s \\

\addlinespace

WebAssembly AOT &
ESP32-CAM / WAMR &
160 MHz &
4 MB PSRAM &
2.000 &
9,591 s \\

\addlinespace

WebAssembly AOT &
ESP32-CAM / WAMR &
240 MHz &
4 MB PSRAM &
2.000 &
6,195 s \\

\addlinespace

WebAssembly desktop &
Windows x64 / Node.js v22.19.0 &
- MHz\textsuperscript{*} &
15,74 GB RAM &
2.000 &
0,041 s \\

\bottomrule
\end{tabularx}

\end{table}

\begin{table}[H]
\centering
\caption{Tempos de inferência e aceleração das configurações avaliadas.}
\label{tab:tempos-inferencia-aceleracao}

\begin{tabular}{
l
r
r
r
r
}
\toprule
\textbf{Configuração} &
\textbf{Média} &
\textbf{Mínimo} &
\textbf{Máximo} &
\textbf{Aceleração} \\
\midrule

Interpretado, 160 MHz &
144,019 s &
143,891 s &
144,436 s &
1,00$\times$ \\

\addlinespace

AOT, 160 MHz &
9,591 s &
9,576 s &
9,613 s &
15,02$\times$ \\

\addlinespace

AOT, 240 MHz &
6,195 s &
6,189 s &
6,202 s &
23,25$\times$ \\

\addlinespace

Desktop &
0,041 s &
0,036 s &
0,170 s &
3479,56$\times$ \\

\addlinespace

\shortstack[l]{TFLM + ESP-NN,\\ESP32-CAM, 240 MHz} &
$\approx 1{,}7$ s &
--- &
--- &
--- \\

\bottomrule
\end{tabular}

\vspace{0.3cm}

\small \textbf{Fonte:} Elaborado pelo autor.

\par\smallskip
\begin{minipage}{\textwidth}
\small \textbf{Nota:} A linha TFLM com ESP-NN apresenta o tempo aproximado informado para a execução do mesmo modelo no ESP32-CAM a 240~MHz. Não estão disponíveis, para essa medição, mínimo, máximo ou estatísticas de dispersão; por isso, esses campos e a aceleração não são preenchidos. A diferença em relação ao WASM AOT caracteriza o desempenho relativo de implementações com kernels e otimizações distintos, sem isolar o custo exclusivo da portabilidade.
\end{minipage}

\end{table}

\begin{table}[H]
\centering
\caption{Acurácia das implementações avaliadas.}
\label{tab:acuracia-configuracoes}

\small
\begin{tabular}{lrrrr}
\toprule
\textbf{Configuração} &
\textbf{Acertos} &
\textbf{Erros} &
\textbf{Sem predição} &
\textbf{Acurácia} \\
\midrule

WASM ESP32-CAM &
1.966 & 29 & 5 & 98,30\% \\

WASM desktop &
1.965 & 30 & 5 & 98,25\% \\

TFLite ESP32-CAM &
1.990 & 9 & 1 & 99,50\% \\

TFLite desktop &
1.990 & 8 & 2 & 99,50\% \\

\bottomrule
\end{tabular}

\vspace{0.3cm}
\small \textbf{Fonte:} Elaborado pelo autor.
\end{table}
